%% GENERATED by python/paper/build.py from results/ -- do not edit by hand.
%% Rebuild:  .venv\Scripts\python.exe python\paper\build.py --only US_CRRA_Ageing
%% Source:   results/shocks/US_shocksCommonX.csv
\begin{table}[!htb]
\centering
\begin{threeparttable}
\caption{Does CRRA matter for the effect of ageing in US -- 2020}
\label{table:US:CRRA:ageing}
\renewcommand{\arraystretch}{1.25}
\begin{tabularx}{.9\textwidth}{YYYYY}
\toprule
\textbf{Scenario} & \textbf{CRRA} ($\rho$) & \textbf{Tax rate} & \textbf{Savings rate} & \textbf{Avg. workweek} \\
\midrule 
 & 0.5 & 14.43\% & 15.91\% & 39.39 \\
Baseline & 1.0 & 14.43\% & 15.37\% & 39.39 \\
 & 2.0 & 14.43\% & 14.85\% & 39.39 \\[.5em]\hline\\[-.75em]
 & 0.5 & 18.76\% & $-1.10$ p.p. & 39.08 \\
Mild ageing & 1.0 & 18.45\% & $-1.00$ p.p. & 39.08 \\
 & 2.0 & 18.01\% & $-0.88$ p.p. & 39.08 \\[.5em]\hline\\[-.75em]
 & 0.5 & 23.75\% & $-2.32$ p.p. & 38.68 \\
Acute ageing & 1.0 & 23.11\% & $-2.08$ p.p. & 38.66 \\
 & 2.0 & 21.99\% & $-1.79$ p.p. & 38.72 \\[.5em]\hline\\[-.75em]
\bottomrule
\end{tabularx}
\begin{tablenotes}
\footnotesize
\item Every $\rho$ is separately calibrated: the tax rate is a target and the workweek is normalised against each $\rho$\textquotesingle s own baseline, so both are common to the baseline rows, while the baseline savings rate is a prediction ($\beta$ is identified by the interest rate) and falls with $\rho$. The savings rate is savings relative to GDP; the baseline rows report its level and every scenario row the change against the baseline at the same $\rho$, in percentage points. Common-$X$ calibration: one leisure parameter $X$ shared across income groups, with the hours unit pinned by targeting the observed average workweek, so relative hours are a prediction rather than a calibration target.
\end{tablenotes}
\end{threeparttable}
\end{table}
